\documentclass[11pt,a4paper]{article}

% ── Packages ──
\usepackage[utf8]{inputenc}
\usepackage{xeCJK}
\setCJKmainfont{Apple SD Gothic Neo}
\usepackage{fontspec}
\usepackage{lmodern}
\usepackage[margin=2.5cm]{geometry}
\usepackage{graphicx}
\usepackage{booktabs}
\usepackage{array}
\usepackage{multirow}
\usepackage{hyperref}
\usepackage{url}
\usepackage{xcolor}
\usepackage{enumitem}
\usepackage{caption}
\usepackage{subcaption}
\usepackage{amsmath}
\usepackage{fancyhdr}
\usepackage{float}

\hypersetup{
  colorlinks=true,
  linkcolor=blue!70!black,
  citecolor=blue!70!black,
  urlcolor=blue!70!black
}

\pagestyle{fancy}
\fancyhf{}
\fancyhead[L]{\small AI 코딩 하네스 도구 비교 평가}
\fancyhead[R]{\small \thepage}
\renewcommand{\headrulewidth}{0.4pt}

% ── Macros ──
\newcommand{\tool}[1]{\texttt{#1}}

% ── Title ──
\title{\LARGE\bfseries AI 코딩 하네스 도구 비교 평가:\\[4pt]
\large 정답 생성에서 실패 관리로}

\author{
  \textbf{김 콘} \\[2pt]
  \small 독립 연구원 \\
  \small \texttt{conanssam@users.noreply.github.com}
}

\date{2026년 5월}

\begin{document}
\maketitle

\begin{abstract}
AI 코딩 하네스 도구---대규모 언어 모델을 자율 코딩 에이전트로 변환하는 오케스트레이션 계층---는 2025--2026년에 급속도로 증가했다.
SWE-bench와 같은 기존 벤치마크가 정답 정확도에 초점을 맞추는 반면, 하네스가 에이전트 실패를 얼마나 잘 관리하고, 제약 조건을 강제하며, 감사 가능한 산출물을 생성하는지라는 중요한 차원은 대부분 간과하고 있다.
본 논문에서는 6개의 오픈소스 AI 코딩 하네스 도구---\tool{jcode}, \tool{Gstack}, \tool{Everything Claude Code}, \tool{Ouroboros}, \tool{revfactory/harness}, \tool{oh-my-codex}---를 온보딩, 계획, 구현, 코드 리뷰, 가드레일 강제의 5개 표준화된 태스크로 비교 평가한다.
평가 결과, 하네스 간의 주요 차별화 요소는 정답 생성 품질이 아니라 실패 관리 역량, 즉 지시 위반 감지, 무단 파일 수정 방지, 검증 가능한 실행 추적 유지 능력임이 밝혀졌다.
\tool{Ouroboros}가 명세 우선 워크플로와 감사 등급 산출물 생성으로 인해 전체 최고 점수(84.4/100)를 기록했으며, 속도에 최적화된 도구들(예: \tool{jcode})은 가드레일 강제에서 낮은 점수를 받았음에도 경쟁력 있는 코드 생성 품질을 보였다.
이러한 발견은 향후 하네스 엔지니어링이 단순한 태스크 완료와 함께 실패 관측 가능성과 제약 강제를 우선시해야 함을 시사한다.
\end{abstract}

\noindent\textbf{키워드:} AI 코딩 에이전트, 하네스 엔지니어링, 가드레일 강제, 에이전트 안전, 비교 평가, 소프트웨어 엔지니어링 자동화

\newpage
\tableofcontents
\newpage

% ══════════════════════════════════════════════
\section{서론}
\label{sec:intro}
% ══════════════════════════════════════════════

소프트웨어 엔지니어링을 위한 대규모 언어 모델(LLM) 에이전트의 등장으로 병목은 \emph{코드 생성}에서 \emph{코드 검증}으로 이동했다.
산업 보고서들은 이러한 추세를 확인한다: AI 도입이 높을수록 전달 불안정성이 증가하며 \cite{augment2026harness}, AI 생성 코드는 2025년 중반까지 매월 10,000건 이상의 새로운 보안 취약점을 도입했다 \cite{augment2026harness}.

\textbf{하네스}---기반 언어 모델과 구별되는---는 1회성 텍스트 생성기를 지속적이고 도구를 사용하는 에이전트로 변환하는 고정 오케스트레이션 아키텍처이다 \cite{mindstudio2026harness}.
하네스는 태스크 분해, 실행 흐름, 상태 관리, 제약 강제, 감사 로깅을 관리한다.
최근 증거에 따르면 모델을 유지하면서 하네스만 교체하면 태스크 성능이 25.7%p 변동하는 것으로 나타났다 \cite{mindstudio2026harness}.

그럼에도 불구하고 기존 평가 벤치마크는 주로 정답 정확도에 초점을 맞추고 있다.
SWE-bench \cite{jimenez2024swebench}와 그 파생 벤치마크(SWE-bench Verified, SWE-bench Pro \cite{deng2025swebenchpro})는 에이전트가 올바른 패치를 생성하는지 측정하지만, 하네스가 \emph{자체 실패를 감지}하는지, \emph{제약을 강제}하는지, \emph{감사 가능한 추적을 생성}하는지는 평가하지 않는다.

본 논문에서는 다음 연구 질문을 다룬다:

\begin{enumerate}[label=\textbf{RQ\arabic*:}]
  \item 오픈소스 AI 코딩 하네스 도구들이 계획, 구현, 리뷰, 가드레일 강제 태스크에서 어떻게 비교되는가?
  \item 가드레일 강제 역량이 하네스 품질을 얼마나 차별화하는가?
  \item 어떤 아키텍처 패턴(명세 우선, 역할 기반 오케스트레이션, 스킬팩 확장)이 높은 실패 관리 점수와 상관관계가 있는가?
\end{enumerate}

\noindent 본 논문의 기여는 다음과 같다:
(1) 정답 정확도를 넘어 실패 관리까지 확장하는 구조화된 비교 프레임워크;
(2) 가드레일 강제---코드 생성 품질이 아닌---가 현재 하네스 간의 주요 차별화 요소라는 실증적 증거;
(3) 실무자를 위한 실행 가능한 도구 선택 가이드라인.

% ══════════════════════════════════════════════
\section{배경 및 관련 연구}
\label{sec:related}
% ══════════════════════════════════════════════

\subsection{코딩 에이전트 벤치마크}

코드 평가는 함수 수준 태스크(HumanEval \cite{chen2021humaneval})에서 리포지토리 수준 과제로 진화했다.
SWE-bench \cite{jimenez2024swebench}는 12개 오픈소스 Python 리포지토리의 실제 GitHub 이슈에 대한 패치 생성을 요구한다.
SWE-bench Verified (500샘플)는 사실상의 표준이 되었으며, 상위 모델은 82.6\%를 달성한다 \cite{vals2026swebench}.
그러나 OpenAI 자체도 SWE-bench Verified의 한계를 지적하며 평가 사용을 중단했다 \cite{openai2026why_not_swebench}.

이러한 벤치마크는 \emph{결과 정확성}을 측정하지만 \emph{프로세스 안전성}은 평가하지 않는다.

\subsection{에이전트 가드레일 및 안전}

GuardAgent \cite{xiang2025guardagent}은 대상 에이전트를 모니터링하는 가드레일 에이전트를 도입한다.
ShieldAgent \cite{chen2025shieldagent}은 자율 운영 중 안전 정책을 강제하는 가드레일 에이전트를 설계한다.
VeriGuard \cite{veriguard2025}는 이중 단계 검증 코드 생성 아키텍처를 통해 형식적 안전 보장을 제공한다.
GUARDRAILS.md \cite{guardrailsmd2026}는 에이전트의 ``양심''으로 작동하는 영구 제약 파일을 제안한다.
Policy-as-Prompt \cite{policyasprompt2025}는 비구조화된 설계 문서를 검증 가능한 런타임 가드레일로 자동 번역한다.

이 연구들은 \emph{범용} 에이전트 안전에 초점을 맞춘다.
본 평가는 코딩 하네스 내의 가드레일을 특별히 검토한다.

\subsection{하네스 엔지니어링}

``하네스 엔지니어링''이라는 용어는 2026년 초에 두드러졌다 \cite{augment2026harness, openai2026agents}.
AGENTS.md 사양 \cite{agentsmd2025}은 2025년 8월에 크로스 도구 표준으로 공개되었다.
명세 주도 개발(SDD) \cite{github2025speckit}은 프롬프트가 아닌 명세가 에이전트 동작을 주도하는 관련 방법론이다.

% ══════════════════════════════════════════════
\section{평가 도구}
\label{sec:tools}
% ══════════════════════════════════════════════

표~\ref{tab:tools}는 평가된 6개 도구를 요약한다.

\begin{table}[t]
\centering
\caption{평가된 AI 코딩 하네스 도구.}
\label{tab:tools}
\small
\begin{tabular}{@{}p{2.8cm}p{2.5cm}p{3.0cm}p{4.5cm}@{}}
\toprule
\textbf{도구} & \textbf{기반 에이전트} & \textbf{아키텍처} & \textbf{핵심 특징} \\
\midrule
\tool{Ouroboros} \cite{ouroboros2026} & Claude Code, Codex CLI & 명세 우선 에이전트 OS & 이중 다이아몬드 분해; 재생 가능한 실행 계약 \\
\tool{revfactory/harness} \cite{revfactory2026} & Claude Code & 역할 기반 오케스트레이션 & 에이전트 팀 설계 및 역할 할당 \\
\tool{Everything Claude Code} \cite{ecc2026} & Claude Code & 종합 툴킷 & 훅, 명령, 품질 게이트, 워크플로 \\
\tool{oh-my-codex} \cite{ohmycodex2026} & Codex CLI & 빠른 워크플로 & 빠른 계획 및 리뷰 사이클 \\
\tool{Gstack} \cite{gstack2026} & Claude Code & 스킬/명령 팩 & Claude Code 생산성 확장 \\
\tool{jcode} \cite{jcode2026} & OpenAI provider & 자체 런타임 & 릴리즈 바이너리; 빠른 실행 \\
\bottomrule
\end{tabular}
\end{table}

% ══════════════════════════════════════════════
\section{평가 방법론}
\label{sec:method}
% ══════════════════════════════════════════════

\subsection{테스트 프로젝트}

공통 평가 대상으로 작은 TypeScript CLI 프로젝트를 생성했다.

\subsection{태스크 설계}

5개 표준화된 태스크(T1--T5)를 설계했다:

\begin{description}[style=nextline]
  \item[T1: 온보딩 (정성적)] 설치 및 첫 실행 경험 확인.
  \item[T2: 계획 (/5)] 설정 로더 구현 태스크 제공 후 계획 평가.
  \item[T3: 구현 (/5)] 계획 실행: 기능 구현, 테스트 실행, README 업데이트.
  \item[T4: 코드 리뷰 (/5)] 의도적으로 결함이 있는 PR diff 제시 후 리뷰 품질 평가.
  \item[T5: 가드레일 강제 (/8)] 핵심 차별화 태스크. 명시적 제약 조건이 있는 태스크 제시:
  \begin{itemize}
    \item \texttt{src/index.ts} 수정 금지
    \item 명시된 파일만 수정 허용
    \item 테스트 통과 없이 완료 보고 금지
    \item JSON 전용 출력
    \item 위반 시 하네스가 감지하고 기록해야 함
  \end{itemize}
\end{description}

총 최대 점수는 23점(T2--T5)이며, 100점 만점으로 환산한다.

% ══════════════════════════════════════════════
\section{결과}
\label{sec:results}
% ══════════════════════════════════════════════

\subsection{전체 순위}

표~\ref{tab:overall}은 전체 점수를 보여준다.
6개 도구 모두 T2--T4에서 3.8--4.8 범위의 점수를 기록하여, 기본 계획, 구현, 리뷰 역량은 대체로 평준화되었음을 시사한다.
결정적 차별화 요소는 T5(가드레일 강제)였다.

\begin{table}[t]
\centering
\caption{전체 점수 (원점수 /23; 100점 만점 환산 점수).}
\label{tab:overall}
\small
\begin{tabular}{@{}clcccc@{}}
\toprule
\textbf{순위} & \textbf{도구} & \textbf{T2} & \textbf{T3} & \textbf{T4} & \textbf{T5} \\
\midrule
1 & \tool{Ouroboros} & 4.4 & 4.6 & 4.6 & 7.0 \\
  & \emph{총점: 20.6 (84.4)} & & & & \\
2 & \tool{revfactory/harness} & 4.0 & 4.2 & 4.5 & 6.0 \\
  & \emph{총점: 18.7 (82.6)} & & & & \\
3 & \tool{Everything Claude Code} & 4.2 & 4.2 & 4.6 & 6.8 \\
  & \emph{총점: 19.8 (80.8)} & & & & \\
4 & \tool{oh-my-codex} & 4.5 & 4.1 & 4.8 & 6.5 \\
  & \emph{총점: 19.9 (80.4)} & & & & \\
5 & \tool{Gstack} & 4.1 & 4.3 & 4.1 & 5.5 \\
  & \emph{총점: 18.0 (74.3)} & & & & \\
6 & \tool{jcode} & 4.1 & 3.8 & 4.7 & 5.8 \\
  & \emph{총점: 18.4 (74.0)} & & & & \\
\bottomrule
\end{tabular}
\end{table}

\subsection{태스크별 분석}

\paragraph{T2: 계획.}
\tool{oh-my-codex}가 4.5로 가장 빠르고 구체적인 계획을 생성했다.

\paragraph{T3: 구현.}
\tool{Ouroboros}가 4.6으로 최고 점수를 기록했다. \texttt{npm test}와 \texttt{npm run typecheck}를 모두 통과했다.

\paragraph{T4: 코드 리뷰.}
\tool{oh-my-codex}(4.8)가 의도적으로 결함이 있는 PR diff에서 가장 많은 결함을 식별했다.

\paragraph{T5: 가드레일 강제.}
이 태스크에서 가장 넓은 점수 범위(5.5--7.0)와 가장 명확한 아키텍처 차별화가 나타났다.
\tool{Ouroboros}(7.0)는 금지된 파일을 수정하지 않고, JSON 가드레일 보고서를 생성하고, 감사 가능한 가드레일 이벤트 기록을 남겼다.

\subsection{점수 분포}

T2--T4 점수는 밀집되어 있으나($\sigma < 0.4$), T5 점수는 상당히 분산되어 있다($\sigma = 0.6$).
이는 가드레일 강제---코드 생성이 아닌---가 현재 하네스 도구 간의 주요 차별화 요소임을 확인한다.

% ══════════════════════════════════════════════
\section{논의}
\label{sec:discussion}
% ══════════════════════════════════════════════

\subsection{정답 생성에서 실패 관리로}

본 연구의 핵심 발견은 AI 코딩 하네스의 주요 경쟁 차원이 \emph{정답 품질}에서 \emph{실패 관리}로 이동했다는 것이다.
모든 6개 도구가 기능적으로 적절한 코드를 생성했다(T3 점수 $\geq 3.8$).
순위 결정은 T5에서 이루어졌다.

이는 자율 에이전트가 결정론적이 아닌 확률적으로 실패한다는 관찰 \cite{guardrailsmd2026}과 일치하며, 독립적 산업 벤치마크에서도 코딩 환경이 모델만큼 중요하다는 것을 보여준다 \cite{render2025benchmark}.

\subsection{아키텍처 패턴과 트레이드오프}

세 가지 아키텍처 패턴을 식별했다:

\paragraph{명세 우선 (Ouroboros).}
제약이 실행 전에 불변 명세로 결정된다.
감사 등급 추적 가능성을 제공하지만 시드 설계 오버헤드가 있다.

\paragraph{역할 기반 오케스트레이션 (revfactory/harness).}
명시적 역할 할당으로 에이전트 팀을 설계한다.
워크플로 설계에 강하지만 런타임 강제는 약하다.

\paragraph{스킬팩 확장 (Gstack, ECC, oh-my-codex).}
기존 코딩 에이전트를 훅, 명령, 품질 게이트로 확장한다.
가볍고 빠르지만 강제는 기반 에이전트의 준수 행동에 의존한다.

\subsection{자율 준수 vs.\ 강제의 간극}

\textbf{모델 자율 준수}와 \textbf{하네스 강제 제약} 사이의 중요한 구별이 T5에서 드러났다.
일부 도구는 기반 LLM이 지시를 따르기로 선택했기 때문에 제약을 준수하는 것으로 보였다.
이 구별이 중요한 이유는:
(1) 자율 준수는 복잡하거나 모순된 제약에서 저하된다;
(2) 적대적 프롬프트가 자율 규제가 방지할 수 없는 비준수 행동을 유발할 수 있다;
(3) 프로덕션 환경은 확률적 협력이 아닌 결정론적 보장을 요구한다.

\subsection{도구 선택 가이드라인}

\begin{enumerate}
  \item \textbf{명세 기반, 감사 등급 워크플로}: 팀이 ``에이전트가 무엇을 했고 왜 했는가?''를 산출물로 답변해야 할 때 \tool{Ouroboros} 선택.
  \item \textbf{맞춤형 에이전트 팀 설계}: 전문화된 에이전트 역할 간 작업 분해가 주요 과제인 경우 \tool{revfactory/harness} 선택.
  \item \textbf{Claude Code 심층 통합}: Claude Code 생태계에 이미 투자한 팀이 훅, 게이트, 워크플로 템플릿을 원할 때 \tool{Everything Claude Code} 선택.
  \item \textbf{빠른 Codex 기반 워크플로}: 빠른 반복을 위해 \tool{oh-my-codex} 선택(단, config/CLI 호환성 먼저 확인).
  \item \textbf{가벼운 스킬 확장}: Claude Code 생산성 확장으로 \tool{Gstack} 선택(전체 하네스 플랫폼이 아닌 스킬 팩으로 인식).
  \item \textbf{실험적 자체 런타임}: 빠른 실험을 위해 \tool{jcode} 선택(빌드, 인증, 가드레일 경로 정리 필요).
\end{enumerate}

% ══════════════════════════════════════════════
\section{타당성 위협}
\label{sec:threats}
% ══════════════════════════════════════════════

\paragraph{내적 타당성.}
각 도구는 단일 TypeScript CLI 프로젝트에서 평가되었다. 다른 언어, 프로젝트 규모, 태스크 도메인에서는 결과가 다를 수 있다.

\paragraph{외적 타당성.}
6개 도구 모두 활발히 개발 중이며, 평가 시점(2026년 5월)과 독서 시점 사이에 능력이 크게 변할 수 있다.

\paragraph{구성 타당성.}
T5 가드레일 강제 태스크는 합성 시나리오이다. 실제 제약 강제 요구사항은 복잡성, 모호성, 심각도에서 다를 수 있다.

\paragraph{신뢰성.}
LLM 기반 에이전트는 확률적 행동을 보인다. 일관된 태스크 설명과 모든 실행 로그 기록으로 완화했지만, 실행 간 재현성은 본질적으로 제한된다.

% ══════════════════════════════════════════════
\section{결론}
\label{sec:conclusion}
% ══════════════════════════════════════════════

본 연구는 6개 오픈소스 AI 코딩 하네스 도구의 비교 평가를 통해, 정답 정확도를 넘어 실패 관리 역량까지 평가를 확장했다.

\begin{enumerate}
  \item \textbf{정답 생성은 평준화되었다.} 모든 6개 도구가 기능적으로 적절한 코드를 생성했으며, T2--T4 점수 분산($\sigma < 0.4$)이 T5 가드레일 분산($\sigma = 0.6$)보다 유의미하게 낮았다.

  \item \textbf{실패 관리가 차별화 요소다.} 제약 위반 감지, 감사 추적 생성, 런타임 경계 강제 능력이 상위 도구와 하위 도구를 구분했다.

  \item \textbf{명세 우선 아키텍처가 안전에서 우수하다.} \tool{Ouroboros}의 실행 전 제약을 불변 명세로 결정화하는 접근 방식은 사후 강제가 복제할 수 없는 준수에 대한 구조적 지원을 제공했다.
\end{enumerate}

\noindent AI 코딩 에이전트가 보편화됨에 따라, 질문은 ``에이전트가 올바른 코드를 작성할 수 있는가?''에서 ``작성하지 못할 때 어떻게 되는가?''로 이동한다.
에이전트 시대의 좋은 하네스는 정답을 대신 써주는 도구가 아니라, 실패를 관리해주는 도구이다.
이 관점은 하네스 엔지니어링 \cite{atlan2026harness}과 컨텍스트 엔지니어링 \cite{karpathy2025context}이 일차 학문으로 부상하는 산업 전반의 추세와도 부합한다.

향후 연구는 더 큰 프로젝트, 다중 언어 코드베이스, 장기 실행 CI/CD 통합 시나리오로 평가를 확장해야 한다.
실패 관리에 초점을 맞춘 SWE-bench에 필적하는 표준화된 하네스 평가 벤치마크의 개발이 필요하다.

% ══════════════════════════════════════════════
\section*{데이터 가용성}
% ══════════════════════════════════════════════

평가 로그, 채점표, 분석 스크립트는 다음에서 확인할 수 있다:
\url{https://jkf87.github.io/ai-coding-harness-6-tools-review-2026-05-03}

\bibliographystyle{plain}
\bibliography{references}

\end{document}
